\subsection{Scenarios, simulation and sensitivity}

\begin{qfmethod}[Historical scenarios, Monte Carlo against historical simulation, and sensitivity analysis for an investment decision: what each answers, how to set one up, and how to read the results]\label{meth:scenario-simulation-analysis}
\qffield{Statement} \emph{Scenario analysis} values an investment under a few discrete states. Damodaran gives four components \citep[pp.~894--897]{investment-valuation-2012}:
\begin{enumerate}
\item the two or three factors the scenarios turn on;
\item the number of scenarios;
\item the outcome in each;
\item the probabilities, if a weighted value is wanted.
\end{enumerate}
\begin{itemize}
\item Best and worst cases set each input to its extreme, which may be infeasible; consistent scenarios respect the links between inputs (pp.~894--895).
\item Scenarios can be ``plausible future narratives'' rather than probable outcomes; their value is the range (p.~896).
\item The weakness is ``garbage in, garbage out'': the scenarios must be realistic and cover the spectrum (p.~897).
\end{itemize}
\emph{Historical scenarios} replay actual episodes on today's position: they are realistic by construction but cannot include what has not happened (\hyperref[meth:historical-simulation]{historical simulation}).

\emph{Simulation} examines continuous risk (pp.~908--910):
\begin{enumerate}
\item choose the few variables that matter;
\item give them distributions from historical, cross-sectional or statistical sources;
\item check correlations among them;
\item run many draws.
\end{enumerate}
\begin{itemize}
\item \emph{Monte Carlo} samples from a fitted model, so it can extend beyond the data but inherits the model's shape (\hyperref[meth:monte-carlo-simulation]{Monte Carlo}).
\item \emph{Historical simulation} resamples actual data, so it keeps fat tails and the joint behaviour of the variables, but only those seen (\hyperref[meth:bootstrap-resampling]{bootstrap}).
\end{itemize}
Simulation gives a full distribution where scenario analysis gives a few points (p.~919).

\emph{Sensitivity analysis} changes one assumption at a time and records the effect on a chosen statistic. Using the same random draws in the base and changed runs (common random numbers) isolates the change from simulation noise (Interpretation).

\qffield{Practice} \emph{A 60/40 portfolio} of the S\&P~500 and 10-year Treasuries, rebalanced yearly and held ten years (\texttt{trading\_notes.numerical.scenario\_simulation.scenario\_simulation\_checks}; data 1928--2025 \citep[sheet ``Returns by year'']{us-asset-returns-1928-2025}; 20,000 paths each).

\emph{Historical scenarios}:
\begin{center}\small
\begin{tabular}{@{}lrl@{}}
\toprule
Episode & Portfolio return & Note \\
\midrule
1931 & $-27.3\%$ & Depression \\
1973--74 & $-20.8\%$ & $-35.2\%$ real after 22.1\% inflation \\
2008 & $-13.9\%$ & Stocks $-36.6\%$, Treasuries $+20.1\%$ \\
2022 & $-18.0\%$ & Stocks $-18.0\%$, Treasuries $-17.8\%$ \\
\bottomrule
\end{tabular}
\end{center}
Treasuries cushioned 2008 but not 2022, so the hedge depends on the source of the shock. A three-scenario view (20 percent recession at $-12$ percent, 60 percent base at $+6$, 20 percent boom at $+18$) has an expected return of 4.8 percent and a one-in-five chance of a double-digit loss.

\emph{Historical simulation against Monte Carlo}, ten-year wealth per 1 invested:
\begin{center}\small
\begin{tabular}{@{}lrrrr@{}}
\toprule
 & 5th pct. & Median & 95th pct. & Loss over 10 years \\
\midrule
Bootstrap of whole years & 1.21 & 2.26 & 3.92 & 1.8\% \\
Normal with the same mean and covariance & 1.22 & 2.23 & 3.97 & 1.4\% \\
\bottomrule
\end{tabular}
\end{center}
\begin{itemize}
\item The two agree on the centre: mean return 9.0 percent, volatility 12.1 percent.
\item They differ in the tail. The one-year 5th percentile is $-13.4$ percent in the data against $-10.8$ percent under the normal, and the chance of losing money over ten years is 1.8 against 1.4 percent.
\end{itemize}

\emph{Sensitivity} of the Monte Carlo result, with common random numbers:
\begin{center}\small
\begin{tabular}{@{}lrr@{}}
\toprule
Change & 5th percentile & Median \\
\midrule
Equity mean 2 points lower & $-0.14$ & $-0.24$ \\
Stock--bond correlation 0.5 (historical 0.02) & $-0.10$ & $-0.03$ \\
80/20 instead of 60/40 & $-0.11$ & $+0.20$ \\
\bottomrule
\end{tabular}
\end{center}
\emph{Interpretation.} The median depends on the return assumption, the tail on the correlation; a move to 80/20 buys median wealth with tail risk. For a client whose constraint is the downside, the correlation assumption, which 2022 showed can change, deserves as much scrutiny as the equity premium.

\qffield{Conditions} Bootstrapping whole years keeps the stock--bond pairing within each year but not dependence across years, and it assumes the 1928--2025 distribution is today's. The normal model assumes constant parameters. Neither includes fees, taxes or rebalancing costs.

\qffield{Pitfalls} Presenting the Monte Carlo median as a forecast. Choosing scenarios that confirm a view. Changing inputs one at a time when they move together, as returns and correlations do in a crisis.

\qffield{See also} \hyperref[meth:historical-simulation]{historical simulation}; \hyperref[meth:monte-carlo-simulation]{Monte Carlo}; \hyperref[meth:sensitivity-scenario-risk-measures]{scenario risk measures}; \hyperref[meth:forecast-model-valuation]{forecast-model scenarios}.

\end{qfmethod}
